\subsection{Questão 73}

\subsubsection{Enunciado}

Um circuito \(LC\) oscila com uma frequência de \(10{,}4\,\text{kHz}\).

\begin{enumerate}
    \item[(a)] Se a capacitância é \(340\,\mu\text{F}\), qual é a indutância?
    
    \item[(b)] Se a corrente máxima é \(7{,}20\,\text{mA}\), qual é a energia total do circuito?
    
    \item[(c)] Qual é a carga máxima do capacitor?
\end{enumerate}


\subsubsection{Resolução}

A frequência angular do circuito é

\[
\omega=2\pi f
=2\pi(10{,}4\times 10^3)
=6{,}53\times 10^4\,\text{rad/s}.
\]

\begin{enumerate}
    \item[(a)] Para um circuito \(LC\),

    \[
    \omega=\frac{1}{\sqrt{LC}}.
    \]

    Portanto,

    \[
    L=\frac{1}{\omega^2 C}.
    \]

    Substituindo os valores,

    \[
    L=
    \frac{1}{(6{,}53\times 10^4)^2(340\times 10^{-6})}
    =
    6{,}89\times 10^{-7}\,\text{H}.
    \]

    \item[(b)] A energia total é igual à energia máxima no indutor:

    \[
    U=\frac{1}{2}LI_{\max}^2.
    \]

    Logo,

    \[
    U=
    \frac{1}{2}(6{,}89\times 10^{-7})(7{,}20\times 10^{-3})^2
    =
    1{,}79\times 10^{-11}\,\text{J}.
    \]

    \item[(c)] Como

    \[
    I_{\max}=\omega Q_{\max},
    \]

    temos

    \[
    Q_{\max}=\frac{I_{\max}}{\omega}
    =
    \frac{7{,}20\times 10^{-3}}{6{,}53\times 10^4}
    =
    1{,}10\times 10^{-7}\,\text{C}.
    \]
\end{enumerate}

\vspace{1cm}
